\documentclass[11pt,a4paper]{article}

\usepackage[utf8]{inputenc}
\usepackage[T1]{fontenc}
\usepackage[spanish]{babel}
\usepackage[margin=2cm, top=2.6cm, bottom=1.8cm, footskip=0.6cm]{geometry}
\usepackage{graphicx}
\graphicspath{{../../01_Propuesta/Logos/}{./}}
\usepackage{fancyhdr}
\usepackage{xcolor}
\usepackage[hidelinks]{hyperref}
\usepackage{enumitem}
\setlist{nosep}
\setlength{\parindent}{0pt}
\setlength{\parskip}{4pt}

\definecolor{verdeosc}{RGB}{26,107,60}
\definecolor{verdemid}{RGB}{46,139,87}

\setlength{\headheight}{40pt}
\pagestyle{fancy}
\fancyhf{}
\fancyhead[L]{%
  \begin{tabular*}{\textwidth}{@{\extracolsep{\fill}}cccc@{}}
    \includegraphics[height=1.0cm]{logo_unisucre.png} &
    \includegraphics[height=1.0cm]{logo_sicvec2026.png} &
    \includegraphics[height=1.0cm]{Logo_insilico.jpeg} &
    \includegraphics[height=1.0cm]{Logo_Bio.png} \\
  \end{tabular*}%
}
\fancyfoot[C]{\footnotesize\itshape SICVEC 2026 contribuye a los Objetivos de Desarrollo Sostenible 12, 13, 15 y 17.}
\renewcommand{\headrulewidth}{1pt}
\renewcommand{\headrule}{\hbox to\headwidth{\color{verdeosc}\leaders\hrule height\headrulewidth\hfill}}

\begin{document}

\vspace{0.1cm}

\noindent
\textbf{Sincelejo, 7 de octubre de 2026}

\vspace{0.15cm}

\noindent
Doctor\\
Wilson Manuel Castro Silupu\\
Profesor Principal y Docente Investigador\\
Universidad Nacional de Frontera (UNF)\\
Sullana, Piura, Perú

\vspace{0.15cm}

\noindent
\textbf{\large Asunto: Difusión de la participación virtual gratuita — SICVEC 2026}

\vspace{0.15cm}

Estimado Dr. Castro:

\vspace{0.1cm}

Como conferencista magistral confirmada del \textbf{Simposio Internacional de Ciencia Verde y Economía Circular (SICVEC 2026)}, quisiéramos pedirle un favor adicional.

El simposio ofrece \textbf{participación virtual gratuita} para público externo e internacional, con certificado de participación, mediante una inscripción en línea sencilla y sin costo para esa modalidad. Nos encantaría que colegas, estudiantes e investigadores de la UNF pudieran acompañarnos, especialmente durante su conferencia. Le agradeceríamos mucho que comparta esta información en su institución, por los canales que considere adecuados.

\vspace{0.1cm}

\noindent
\textbf{INFORMACIÓN PARA COMPARTIR}

\begin{tabular}{ll}
\textbf{Evento:} & Simposio Internacional de Ciencia Verde y Economía Circular (SICVEC 2026) \\
\textbf{Fechas:} & 19 y 20 de octubre de 2026 (lunes y martes) \\
\textbf{Modalidad:} & Híbrida — presencial en Sincelejo, Colombia / virtual gratuita (externos e internacionales) \\
\textbf{Inscripción virtual (gratuita):} & \href{https://sicvec2026.eu.pythonanywhere.com/inscripcion}{sicvec2026.eu.pythonanywhere.com/inscripcion} \\
\textbf{Certificado:} & de participación, para quienes se inscriban y asistan \\
\textbf{Programa y conferencistas:} & \href{https://sicvec2026.eu.pythonanywhere.com/programa}{sicvec2026.eu.pythonanywhere.com/programa} \\
\end{tabular}

\vspace{0.1cm}

Muchas gracias de antemano por su apoyo en la difusión, y por acompañarnos con su conferencia en SICVEC 2026.

\vspace{0.15cm}

Atentamente,

\vspace{0.15cm}

\textbf{Dr. Aldo F. Combariza}\\
Coordinador General, SICVEC 2026\\
Grupo de Investigación IN SILICO · Departamento de Biología y Química\\
Facultad de Educación y Ciencias · Universidad de Sucre\\
\href{mailto:aldo.combariza@unisucre.edu.co}{aldo.combariza@unisucre.edu.co}

\end{document}
